\begin{blocksection}
\question Based on the Link class that is below, write the scheme expression represented by the Link class.  
Assume that this is the Link class in Python that we will be using:
\begin{lstlisting}
class Link:
    '''A Scheme list is a Link in which rest is a Link or nil.'''
    empty = ()
    def __init__(self, first, rest=empty):
        self.first = first
        self.rest = rest

    # There are also __str__, __repr__, and map methods, omitted here.

nil = Link.empty
>>> Link('*', Link(5, Link(Link('-', Link(10, Link(2, nil))), nil)))
\end{lstlisting}
\begin{solution}
(* 5 (- 10 2))
\end{solution}

\begin{lstlisting}
>>> Link('or', Link(Link('>', Link(5, Link(2, nil))), Link(Link('/', Link(1, Link(2, nil))), nil)))
\end{lstlisting}
\begin{solution}
\verb! (or (> 5 2) (/ 1 2))!
\end{solution}

\begin{lstlisting}
    >>> Link('+', Link(1, Link(Link('*', Link(Link('-', Link(5, Link(2, nil))), Link(Link('+', Link(3, Link(1, nil))), nil))), Link(5, nil))))
\end{lstlisting}

\begin{solution}
\verb! (+ 1 (* (- 5 2) (+ 3 1)) 5)!
\end{solution}
    
\end{blocksection}



\begin{guide}
    \begin{blocksection}
        \textbf{Teaching Tips}
        \begin{itemize}
        \item Explain the Link class thoroughly and if you have time, explain what its connection is to interpreters and the scheme project. 
        \end{itemize}
    \end{blocksection}
    \end{guide}